%!TEX program = xelatex
\def\numeromodulo{1}
\def\numerolista{3}
\documentclass[12pt]{exam}

\def\ano{2023}
\def\semestre{1}
\def\disciplina{Introdução à \'Algebra Linear}
\def\turma{3}

\usepackage{caption}
\usepackage{amssymb}
\usepackage{amsmath,amsfonts,amsthm,amstext}
\usepackage[utf8]{inputenc} % Required for inputting international characters
\usepackage[T1]{fontenc}
\usepackage[brazil]{babel}
\usepackage{graphicx}
\graphicspath{{../figuras_geogebra/}{../Pictures/}}
\usepackage{enumitem}
\usepackage{multicol}
\usepackage{answers}
\usepackage[svgnames]{xcolor}
\usepackage{tikz}
\usepackage{ifthen}
\usetikzlibrary{lindenmayersystems}
\usetikzlibrary[shadings]
\usepackage{tikz-3dplot}
\usepackage{tcolorbox}
%-------------amatrix
% Augmented matrix.  Usage (note the argument does not count the aug col):
% \begin{amatrix}{2}
%   1  2  3 \\  4  5  6
% \end{amatrix}
\newenvironment{amatrix}[1]{%
  \left(\begin{array}{@{}*{#1}{r}|r@{}}
}{%
  \end{array}\right)
}

\Newassociation{solucao}{Solution}{ans}
\newtheorem{exercicio}{}

\setlength{\topmargin}{-1.0in}
\setlength{\oddsidemargin}{0in}
\setlength{\textheight}{10.1in}
\setlength{\textwidth}{6.5in}
\setlength{\baselineskip}{12mm}

\newcounter{exercicios}
\setcounter{exercicios}{0}
\newcommand{\questao}{
\addtocounter{exercicios}{1}
\noindent{\bf Exerc{\'\i}cio \Roman{exercicios}: }}

\newcommand{\resp}[1]{
\noindent{\bf Exerc{\'\i}cio #1: }}

\extraheadheight{0.7in}
\firstpageheadrule
\runningheadrule
\lhead{
        \begin{minipage}[c]{1.7cm}
        \includegraphics[width=1.7cm]{unb.pdf}
        \end{minipage}%
        \hspace{0pt}
        \begin{minipage}[c]{4in}
          {Universidade de Bras{\'\i}lia} --
          {Departamento de Matem{\'a}tica}
\end{minipage}
\vspace*{-0.8cm}
}
% \chead{Universidade de Bras{\'\i}lia - Departamento de Matem{\'a}tica}
% \rhead{}
% \vspace*{-2cm}

\extrafootheight{.5in}
\footrule
\lfoot{\disciplina\ - \semestre$^o$/\ano\ - Módulo \numeromodulo}
\cfoot{}
\rfoot{P\'agina \thepage\ de \numpages}
\newcommand{\nsub}{\varsubsetneq}
\newcommand{\vaz}{\emptyset}
\newcommand{\im}{{\rm Im\,}}
\newcommand{\sub}{\subseteq}
\renewcommand{\sin}{{\rm sen\,}}
\newcommand{\n}{\mathbb{N}}
\newcommand{\z}{\mathbb{Z}}
\newcommand{\rac}{\mathbb{Q}}
\newcommand{\real}{\mathbb{R}}
\newcommand{\complex}{\mathbb{C}}
\newcommand{\cp}[1]{\mathbb{#1}}
\newcommand{\ch}{\mbox{\textrm{car\,}}\nobreak}
\newcommand{\vesp}[1]{\vspace{ #1  cm}}
\newcommand{\compcent}[1]{\vcenter{\hbox{$#1\circ$}}}
\newcommand{\comp}{\mathbin{\mathchoice
{\compcent\scriptstyle}{\compcent\scriptstyle}
{\compcent\scriptscriptstyle}{\compcent\scriptscriptstyle}}}

\newcommand{\integer}{\mathbb{Z}}
\begin{document}

  \Opensolutionfile{ans}[ans1]
  \begin{center}
    {\Large\bf \disciplina\ - Turma \turma\ -- \semestre$^{o}$/\ano} \\ \vspace{9pt} {\large\bf
        $\numerolista^a$ Lista de Exerc{\'\i}cios -- Módulo \numeromodulo}\\ \vspace{9pt} Prof. Jos{\'e} Ant{\^o}nio O. Freitas
  \end{center}
  \hrule

\vesp{.6}

Nos exerc{\'\i}cios \ref{matrizinicio} \`a \ref{matrizfim}, mostre que $B$ \'e a inversa de $A$.
\begin{exercicio}\label{matrizinicio}
  \[
    A =\begin{bmatrix}
      1 & -1\\
      2 & 3
    \end{bmatrix}, B =\begin{bmatrix}
      3/5 & 1/5\\
      -2/5 & 1/5
    \end{bmatrix}
  \]
\end{exercicio}

\begin{exercicio}
  \[
    A =\begin{bmatrix}
      -2 & 2 & 3\\
      1 & -1 & 0\\
      0 & 1 & 4
    \end{bmatrix}, B =\begin{bmatrix}
      -4/3 & -5/3 & 1\\
      -4/3 & -8/3 & 1\\
      1/3 & 2/3 & 0
    \end{bmatrix}
  \]
\end{exercicio}

\begin{exercicio}\label{matrizfim}
  \[
    A =\begin{bmatrix}
      2 & -17 & 11\\
      -1 & 11 & -7\\
      0 & 3 & -2
    \end{bmatrix}, B =\begin{bmatrix}
      1 & 1 & 2\\
      2 & 4 & -3\\
      3 & 6 & -5
    \end{bmatrix}
  \]
\end{exercicio}

Nos exerc{\'\i}cios \ref{matrizinversainicio} \`a \ref{matrizinversafim}, encontre a inversa da matriz dada, se existir.

\begin{exercicio}\label{matrizinversainicio}
  $
    A =\begin{bmatrix}
        1 & 2\\
        3 & 7
    \end{bmatrix}
  $
  \begin{solucao}
    $A^{-1} =\begin{bmatrix}
      7 & -2\\
      -3 & 1
    \end{bmatrix}$
  \end{solucao}
\end{exercicio}

\begin{exercicio}
  $
    A =\begin{bmatrix}
        -7 & 33\\
        4 & 19
    \end{bmatrix}
  $
  \begin{solucao}
    $A^{-1} =\begin{bmatrix}
      -19 & -33\\
      -4 & -7
    \end{bmatrix}$
  \end{solucao}
\end{exercicio}

\begin{exercicio}
  $
    A =\begin{bmatrix}
        2 & 4\\
        4 & 8
    \end{bmatrix}
  $
  \begin{solucao}
    N\~ao possui inversa.
  \end{solucao}
\end{exercicio}

\begin{exercicio}
  $
    A =\begin{bmatrix}
        1 & 1 & 1\\
        3 & 5 & 4\\
        3 & 6 & 5
      \end{bmatrix}
    $
  \begin{solucao}
    $
      A^{-1} =\begin{bmatrix}
        1 & 1 & -1\\
        -3 & 2 & -1\\
        3 & -3 & 2
      \end{bmatrix}
    $
  \end{solucao}
\end{exercicio}

\begin{exercicio}
  $
    A =\begin{bmatrix}
        1 & 2 & -1\\
        3 & 7 & -10\\
        7 & 16 & -21
    \end{bmatrix}
  $
  \begin{solucao}
    N\~ao tem inversa.
  \end{solucao}
\end{exercicio}

\begin{exercicio}
  $
    A =\begin{bmatrix}
        -8 & 0 & 0 & 0\\
        0 & 1 & 0 & 0\\
        0 & 0 & 0 & 0\\
        0 & 0 & 0 & -5
    \end{bmatrix}
  $
  \begin{solucao}
   N\~ao existe inversa.
  \end{solucao}
\end{exercicio}

\begin{exercicio}\label{matrizinversafim}
  $
    A =\begin{bmatrix}
        1 & 1 & 2\\
        3 & 1 & 0\\
        -2 & 0 & 3
    \end{bmatrix}
  $
  \begin{solucao}
    $
      A^{-1} =\begin{bmatrix}
        -3/2 & 3/2 & 1\\
        9/2 & -7/2 & -3\\
        -1 & 1 & 1
      \end{bmatrix}
    $
  \end{solucao}
\end{exercicio}

\newpage
\Closesolutionfile{ans}
\hrule
\begin{center}
{\large\bf RESPOSTAS}
\end{center}
\hrule
\input{ans1}

\end{document}
